\par\noindent\SourceLeaf{56}%
This series, multiplied by the prefixed factor, can be resolved using the known formula

\[
\sin\varphi\cos\lambda\varphi=\frac{1}{2}\sin(\lambda+1)\varphi-\frac{1}{2}\sin(\lambda-1)\varphi,
\]

\noindent into terms arranged by the sines of multiples. In this way, every term except the initial and the first is split into two, one taking the multiple immediately below, the other the multiple immediately above, the one whose cosine the term being split contains. For example, the term

\[
\nu n . \sin\varphi\,A^{(2)}\cos2\varphi
\]

\noindent gives

\[
-\nu\frac{nA^{(2)}}{2}\sin\varphi+\nu\frac{nA^{(2)}}{2}\sin3\varphi;
\]

\noindent Of these, the first, when added to $\nu nA\sin\varphi$, which arises from the initial term, gives the initial term of the expanded series

\[
=\left(\nu nA-\nu\frac{nA^{(2)}}{2}\right)\sin\varphi.
\]

In general, suppose I seek the coefficient that $\sin(h+1)\varphi$ acquires in the expanded series. In the given series, the terms $\nu nA^{(h)}\sin\varphi\cos h\varphi$ (the first part of its expansion being

\[
+\nu\frac{nA^{(h)}}{2}\sin(h+1)\varphi)
\]

\noindent and $\nu n.A^{(h+2)}\sin\varphi\cos(h+2)\varphi$ (the latter part of its expansion being

\[
-\nu\frac{nA^{(h+2)}}{2}\sin(h+1)\varphi),
\]

\noindent are the only ones from which terms containing the required multiple of $\varphi$ can arise. Thus the complete coefficient of $\sin(h+1)\varphi$ in this series is

\[
=\nu\frac{n.A^{(h)}}{2}-\nu\frac{nA^{(h+2)}}{2}.
\]

\Needspace{8\baselineskip}
$\beta$) On the other side, however, we shall have

\[
(1+n\cos\varphi)\left(A^{(1)}\sin\varphi+2A^{(2)}\sin2\varphi+.\,.\,.+hA^{(h)}\sin h\varphi+.\,.\,.\,.\,.\right)
\]

\noindent and expanding this expression yields two series. One, namely

\[
A^{(1)}\sin\varphi+2A^{(2)}\sin2\varphi+.\,.\,.\,.+(h+1)A^{(h+1)}\sin(h+1)\varphi+.\,.\,.\,.
\]

\noindent follows at once.
